\chapter{Trabalho, Produtividade e Quem Fica com o Ganho}
\label{ch:trabalho}

\ficcao O relatório da Lavandaria chega com uma boa notícia: o rascunho de resposta passou de vinte minutos para quatro. O gerente prepara o anúncio. A contabilista pede as linhas seguintes. Oito minutos de revisão. Dez de correção. O anúncio espera; a soma dá vinte e dois.

O episódio é inventado, os tempos são construídos e a conta pode ser confirmada. A pergunta que abre é real: quando dizemos que a IA aumentou produtividade, o que ficou mais rápido, quem fez o resto e que resultado foi aceite? A economia da tecnologia começa nesta diferença entre produzir uma saída e concluir trabalho.

Vamos partir da tarefa individual e aumentar a escala: processo, organização e economia. Em cada passagem, será preciso acrescentar evidência. Um ganho no primeiro nível não transporta sozinho a conclusão para os seguintes.

\section{Exposição não é despedimento}
Uma profissão é um conjunto de tarefas, relações, responsabilidades e condições de trabalho. Um índice de exposição estima em que medida determinadas tarefas podem ser afetadas por uma tecnologia. Não observa automaticamente substituição de pessoas.

O índice da OIT de 2025 combina informação sobre tarefas, apreciação humana e apoio de modelos para estimar exposição ocupacional. Aponta para cerca de um quarto do emprego mundial em ocupações com algum grau de exposição, com diferenças importantes entre grupos e países~\citep{ilo2025}. Isto não equivale a «um quarto dos empregos vai desaparecer». Entre capacidade técnica e emprego existem adoção, custos, procura, reorganização, normas e decisões políticas.

Uma organização pode automatizar parte de uma função e aumentar produção; pode redistribuir trabalho; pode reduzir contratação; pode criar novas tarefas. São mecanismos possíveis, não efeitos garantidos. Precisamos de observar quais ocorreram e em que condições.

Pensa numa profissão como um cesto de atividades. Algumas são repetíveis, outras exigem negociação, presença, responsabilidade ou conhecimento de situações particulares. Automatizar a redação de uma minuta não demonstra que desapareceu a necessidade de interpretar o pedido, responder por uma decisão ou acompanhar a sua execução.

O próprio cesto pode mudar. Se escrever ficou mais barato, podem passar a pedir-se mais versões ou mais documentos. Se verificar continua caro, o trabalho pode deslocar-se para revisão. Para saber o efeito sobre pessoas, precisamos de observar tarefas, horas, contratação, salários, qualidade e condições de trabalho. Uma lista de capacidades técnicas responde apenas a uma parte da pergunta.

\section{Resultados que parecem contraditórios}
No artigo final de Brynjolfsson, Li e Raymond, a introdução faseada de um assistente em apoio ao cliente foi estudada com dados de 5172 agentes. A produtividade, medida por questões resolvidas por hora, aumentou em média 15\%, com forte heterogeneidade~\citep{brynjolfsson2025}. O desenho inclui comparação temporal e entre grupos sob hipóteses de identificação; não foi uma experiência que distribuiu aleatoriamente todos os trabalhadores do mundo.

Num ensaio da METR sobre ferramentas do início de 2025, 16 programadores experientes trabalharam em 246 tarefas dos seus próprios projetos maduros. As tarefas foram atribuídas a condições com ou sem IA; o tempo de conclusão aumentou 19\% quando o uso era permitido~\citep{metr2025}. É um retrato daquela população, ferramentas e período, não uma estimativa eterna de programação com IA.

\begin{table}[htbp]
\small
\begin{tabularx}{\textwidth}{@{}lXX@{}}
\toprule
Aspeto & Apoio ao cliente & Programação experiente \\
\midrule
Tarefa & Resolver questões de clientes & Alterar projetos conhecidos \\
Medida & Questões por hora & Tempo para concluir tarefas \\
Desenho & Adoção faseada, análise econométrica & Atribuição aleatória por tarefa \\
Resultado citado & Ganho médio com heterogeneidade & Aumento de tempo nas condições testadas \\
Limite & Uma operação e um assistente & Pequena população, ferramentas de 2025 \\
\bottomrule
\end{tabularx}
\caption{Resultados diferentes não respondem necessariamente à mesma pergunta.}
\end{table}

Não precisamos de escolher o estudo de que gostamos mais. Precisamos de perguntar se o nosso trabalho se aproxima das tarefas, do nível de experiência e do processo de revisão avaliados. A opinião do utilizador sobre rapidez também pode divergir do tempo observado.

Os dois desenhos ajudam ainda a formular hipóteses diferentes. Um assistente pode disponibilizar conhecimento útil a quem ainda está a aprender um procedimento. Um especialista pode gastar mais tempo a rever propostas para um projeto que conhece profundamente. Estas são explicações plausíveis a confrontar com os dados de cada contexto, não uma regra segundo a qual principiantes ganham sempre e especialistas perdem sempre.

É tentador escolher um herói e um vilão. Os estudos não são equipas adversárias num campeonato sobre gostar de IA. São instrumentos para reduzir incerteza em perguntas delimitadas.

\section{A conta que o rascunho deixou fora}
Voltemos aos tempos fictícios da abertura. Admitimos etapas seriais, sem sobreposição, e uma tarefa final com o mesmo critério de qualidade. O processo sem IA demora vinte minutos. O processo com IA inclui quatro de geração, oito de revisão e dez de retrabalho. A soma é vinte e dois: o tempo aumentou 10\%.

Produção por unidade de tempo não é a mesma percentagem com o sinal trocado. Com uma tarefa a cada vinte minutos, seriam três por hora neste modelo simplificado; com uma a cada vinte e dois, seriam cerca de $2.73$. A redução é aproximadamente 9.1\%. A diferença nasce do denominador: numa medida dividimos por vinte, na outra comparamos os respetivos inversos.

Agora construímos uma alternativa que realmente termina em dez minutos. O tempo caiu 50\%; mantendo tarefa e qualidade, a produção por hora duplicou, um aumento de 100\%. Dizer apenas «ganho de cinquenta por cento» deixa em aberto qual destas grandezas foi medida.

\textbf{Aprofundamento opcional.} Se $t_0$ é o tempo inicial e $t_1$ o tempo novo para a mesma unidade de trabalho aceite, a redução relativa do tempo é $1-t_1/t_0$. A variação relativa da produção por tempo é $t_0/t_1-1$. O módulo \texttt{exemplos.laboratorio.decisoes} executa as duas contas. Os valores são sintéticos; não resultam de participantes ou de um ensaio de produtividade real.

Se as etapas podem decorrer em paralelo, a simples soma deixa de representar o tempo de calendário. Se uma geração demora dois minutos enquanto a pessoa faz outra tarefa, esse intervalo também não equivale automaticamente a dois minutos de trabalho humano. Um relatório útil distingue tempo de parede, tempo de pessoas e consumo computacional. Misturá-los pode produzir ganhos impossíveis de localizar na agenda.

\section{O contrafactual: o que teria acontecido sem a ferramenta?}
Comparar esta semana com a anterior pode confundir efeito da IA com sazonalidade, mudança de equipa ou tarefas mais fáceis. Uma experiência aleatorizada, quando adequada e viável, ajuda a separar efeitos. Estudos observacionais podem ser informativos com desenho e hipóteses explícitos. Nenhum método dispensa qualidade da medição.

Define o resultado antes de recolher: tarefa aceite, defeitos posteriores, tempo total, intervenção humana, custo e satisfação. Linhas de código, número de mensagens ou palavras produzidas podem aumentar sem que haja mais valor. Uma empresa que duplicou a produção de relatórios que ninguém lê fez crescer uma coluna de Excel, não necessariamente a produtividade.

Conta o trabalho invisível: formular pedidos, preparar dados, verificar, recuperar de falhas e manter integrações. Se o tempo do revisor fica fora da medida, o ganho pode ter sido apenas transferido para outra pessoa.

Um desenho pequeno pode começar por fixar a unidade de aceitação. Para extrair campos de mensagens sintéticas, por exemplo, define quais os campos obrigatórios, o que conta como erro e como se resolvem ambiguidades. Regista tempo até aceitação, incluindo tentativas rejeitadas. Se uma condição produz mais saídas erradas, comparar apenas as que passaram seleciona os sucessos e esconde o custo das restantes.

Há ainda efeitos de ordem. Quem faz a mesma tarefa duas vezes pode aprender com a primeira, mesmo sem IA. E pessoas que trabalham juntas podem partilhar uma solução entre condições. Distribuir casos equivalentes, variar a ordem quando apropriado e registar dependências ajuda; não cria automaticamente uma experiência capaz de estimar todos os efeitos. O desenho deve conservar as suas limitações.

Para uma organização, convém definir também o que faria mudar a decisão. Se não há ganho líquido, melhora-se o processo, restringe-se a tarefa ou abandona-se a aplicação. Continuar porque o tempo de redação caiu pode ser racional para um objetivo de comunicação específico; não deve ser apresentado como prova de uma redução no custo total.

\section{Uma parte cinco vezes mais rápida não acelera tudo cinco vezes}
Suponhamos, noutro cenário sintético, que produzir o rascunho ocupa 20\% de um processo e o restante ocupa 80\%. Se o rascunho ficar cinco vezes mais rápido e tudo o resto permanecer igual, a duração passa de uma unidade para $0.8+0.2/5=0.84$. A redução total é 16\%, não 80\%. Mesmo eliminando inteiramente essa etapa, continuaríamos com as outras quatro quintas partes.

O pressuposto de «tudo o resto igual» é o ponto de partida da conta, não uma previsão. Na prática, a revisão pode aumentar, diminuir ou mudar. A conta serve para perguntar onde está o obstáculo. Se o maior tempo está à espera de uma aprovação, gerar mais depressa não resolve essa espera. Se está a procurar informação dispersa, a recuperação pode ser mais valiosa do que a redação.

Esta forma de abrir o processo também permite identificar complementaridades. Formatos mais claros podem reduzir integração; dados consistentes podem tornar uma resposta verificável; uma autorização bem definida pode evitar retrabalho. A melhoria pode pertencer ao desenho conjunto e ser mal atribuída a um único modelo. Separar componentes na avaliação ajuda a descobrir o que merece investimento.

\section{Distribuição, competências e entrada na profissão}
Mesmo com um ganho médio, pessoas e organizações podem beneficiar de forma desigual. Quem controla dados, distribuição, equipamento e decisão de adoção pode capturar uma parte diferente do valor de quem executa o trabalho. Salários e emprego não são resultados automáticos de uma melhoria técnica.

A aprendizagem profissional também muda se forem removidas tarefas que serviam de treino a principiantes. É uma hipótese organizacional a avaliar: como se adquire experiência quando o rascunho inicial é sempre produzido por uma máquina? Podemos criar revisão acompanhada, explicar escolhas e reservar tarefas sem assistência. O objetivo é medir competência futura, além da entrega de hoje.

Formação, acessibilidade e participação dos trabalhadores são escolhas de governação. Não há fundamento para prometer segurança de emprego a quem aprende meia dúzia de prompts. Há fundamento para testar competências e tornar transparentes os efeitos de uma mudança de processo.

Uma pergunta distributiva acompanha cada ganho: quem dispõe do tempo libertado? Uma equipa pode usá-lo para reduzir sobrecarga, atender mais pessoas ou aumentar exigências. A escolha não vem codificada na taxa de tokens por segundo. É uma decisão organizacional que deve aparecer na análise de benefícios.

Para quem começa uma profissão, interessa preservar situações em que se constrói julgamento. Pode ser necessário criar exercícios, revisão acompanhada e tarefas graduais que antes apareciam no trabalho corrente. A proposta é uma opção de formação a avaliar, não uma prova de que toda a automatização destrói aprendizagem. O capítulo~\ref{ch:relacoes} separa desempenho assistido de competência transferível.

\section{Do ganho de uma tarefa à economia}
Efeitos macroeconómicos dependem de difusão, investimento complementar, procura, concorrência e reorganização. Um estudo de uma tarefa não estima diretamente o PIB. Uma projeção deve identificar cenários de adoção e produtividade, horizonte e incerteza.

Também pode haver criação de procura: se uma atividade fica mais barata, faz-se mais dela. O benefício e o custo total podem ambos crescer. Esta possibilidade reaparece no capítulo de energia. Não concluas que eficiência unitária reduz automaticamente consumo ou emprego agregado.

A cadeia de inferência deve, por isso, ficar visível. Primeiro observamos um resultado numa tarefa. Depois investigamos se permanece quando muitas pessoas e processos mudam. Por fim, perguntamos como difusão, procura e distribuição alteram os efeitos. Cada salto pede dados novos. Multiplicar uma percentagem de laboratório por todo o emprego de um país não resolve as etapas que ficaram pelo caminho.

O relatório fictício da Lavandaria termina, desta vez, com uma alteração concreta: a equipa vai medir respostas aceites e tempo completo, e não apenas rascunhos gerados. Talvez a ferramenta venha a compensar. A investigação deixou de depender de uma impressão e passou a ter uma condição de sucesso.

\begin{promptlabbox}{Medir a tarefa aceite, incluindo a revisão}
\textbf{Objetivo e contexto:} preparar uma comparação pequena de trabalho fictício ou autorizado. Escolhe uma tarefa repetível, como extrair campos de mensagens sintéticas.

\textbf{Procedimento:} executa \texttt{python3 -m exemplos.laboratorio.decisoes} para confirmar as contas sintéticas de vinte para dez e de vinte para vinte e dois minutos. Para uma extensão com pessoas, define antes aceitação, erros e tempo total; distribui tarefas de dificuldade semelhante entre condições; inclui verificação e correções; regista experiência prévia e versão das ferramentas. Se participares em ambas, trata aprendizagem e ordem como limitações.

\textbf{Resultado esperado e validação:} as contas distinguem redução de tempo e aumento da produção por tempo. O relatório de uma eventual extensão mostra tempo até aceitação e defeitos. Não há resultados com participantes humanos nesta edição; essa extensão é um protocolo proposto.

\textbf{Falhas possíveis:} escolher tarefas depois de ver o desempenho, contar só geração ou transformar exposição em previsão de desemprego.
\end{promptlabbox}

Uma tarefa aceite é um critério adequado para muitas operações de escritório. Na saúde, na ciência ou numa fábrica, é preciso abrir ainda outra porta: que ligação existe entre acertar uma previsão e melhorar o que acontece no mundo? É essa passagem que vamos investigar a seguir.
